%========================================
% MatinBook - Stage 12: Bibliography Test
% Test: Biblatex + Citations + References + Print Bibliography
% Purpose: Validate bibliography management system
%========================================

\makeatletter
\def\input@path{{../}}
\makeatother

\documentclass{matinbook}

% Add bibliography resource file
\addbibresource{../references.bib}

\title{آزمون سیستم منابع و کتابنامه}
\author{تیم توسعه MatinBook}
\date{تابستان ۱۴۰۵}

\begin{document}

%----------------------------------------
% Title Page
%----------------------------------------
\maketitle

%----------------------------------------
% Chapter 1: Introduction to Citations
%----------------------------------------
\chapter{مقدمه‌ای بر استناددهی}

\section{اهمیت استناددهی}

استناددهی\index{استناددهی} یکی از ارکان اساسی
نگارش علمی است. هر کتاب علمی باید منابع مورد استفاده
خود را به طور دقیق مشخص کند. \lr{MatinBook}
از بسته \lr{biblatex} با موتور \lr{biber}
برای مدیریت منابع استفاده می‌کند.

\section{شیوه استناد}

در این کتاب از شیوه \textbf{عددی} (\lr{numeric})
برای استناددهی استفاده می‌شود. هر منبع با یک شماره
مشخص می‌شود و در بخش کتابنامه به ترتیب استناد
فهرست می‌شود.

%----------------------------------------
% Chapter 2: Citation Examples
%----------------------------------------
\chapter{نمونه‌های استناد}

\section{کتاب‌ها}

کتاب‌های مرجع در زمینه علوم کامپیوتر:

\begin{itemize}
    \item \textbf{هنر برنامه‌نویسی:}
    دونالد کنوث \cite{knuth1984} مجموعه‌ای بی‌نظیر
    درباره الگوریتم‌ها و ساختمان داده‌ها نوشته است.
    
    \item \textbf{راهنمای لاتک:}
    لسلی لمپورت \cite{lamport1994} سیستم \LaTeX{}
    را برای حروف‌چینی متون علمی معرفی کرد.
    
    \item \textbf{مقدمه‌ای بر الگوریتم‌ها:}
    کتاب \lr{CLRS} \cite{cormen2009} مرجع استاندارد
    الگوریتم‌ها در دانشگاه‌های سراسر جهان است.
\end{itemize}

\section{مقالات علمی}

\begin{itemize}
    \item \textbf{نسبیت خاص:}
    آلبرت اینشتین در سال ۱۹۰۵ مقاله معروف خود
    درباره الکترودینامیک اجسام متحرک را منتشر کرد
    \cite{einstein1905}. این مقاله پایه‌گذار
    نظریه نسبیت خاص بود.
\end{itemize}

\section{منابع آنلاین}

\begin{itemize}
    \item \textbf{پروژه لاتک:}
    وبسایت رسمی پروژه \LaTeX{} \cite{latex-project}
    منبع معتبری برای دریافت اطلاعات به‌روز است.
\end{itemize}

\section{استنادهای چندگانه}

مطالعه هم‌زمان چند منبع برای درک عمیق یک موضوع
ضروری است. برای مثال، برای یادگیری الگوریتم‌ها
می‌توان به آثار کنوث \cite{knuth1984} و کورمن
\cite{cormen2009} مراجعه کرد.

همچنین برای آشنایی با \LaTeX{}، منابع
\cite{lamport1994,latex-project} توصیه می‌شوند.

%----------------------------------------
% Chapter 3: Citation Commands
%----------------------------------------
\chapter{دستورات استناد}

\section{دستورات پایه}

\lr{MatinBook} از تمام دستورات استاندارد \lr{biblatex}
پشتیبانی می‌کند:

\begin{table}[h]
\centering
\caption{دستورات استناد در \lr{biblatex}}
\label{tab:citation-commands}
\begin{tabular}{|c|C{5cm}|c|}
\hline
\textbf{دستور} & \textbf{توضیحات} & \textbf{مثال} \\
\hline
\lr{\textbackslash cite} & استناد ساده & \cite{knuth1984} \\
\lr{\textbackslash parencite} & استناد پرانتزی & \parencite{lamport1994} \\
\lr{\textbackslash textcite} & استناد متنی & \textcite{cormen2009} \\
\lr{\textbackslash autocite} & استناد خودکار & \autocite{einstein1905} \\
\hline
\end{tabular}
\end{table}

\section{تنظیمات کتابنامه}

تنظیمات \lr{biblatex} در \lr{MatinBook}:

\begin{itemize}
    \item \textbf{موتور:} \lr{biber}
    \item \textbf{شیوه:} \lr{numeric} (عددی)
    \item \textbf{ترتیب:} \lr{none} (به ترتیب استناد)
    \item \textbf{فایل منبع:} \texttt{references.bib}
\end{itemize}

\section{ارجاع به منابع بدون استناد}

می‌توان منابعی را که مستقیماً به آن‌ها استناد نشده
اما مطالعه شده‌اند، با \texttt{\textbackslash nocite}
در کتابنامه گنجاند:

% Uncomment to include all references even if not cited:
% \nocite{*}

%----------------------------------------
% Chapter 4: Bibliography Display
%----------------------------------------
\chapter{نمایش کتابنامه}

\section{فرمت خروجی}

کتابنامه در \lr{MatinBook} با عنوان «منابع و مراجع»
و به صورت راست‌به‌چپ نمایش داده می‌شود.
محیط \lr{latin} برای حفظ جهت صحیح متون انگلیسی
استفاده می‌شود.

\section{فهرست منابع}

در ادامه، تمام منابع استناد شده در این کتاب
نمایش داده می‌شوند:

\begin{latin}
\printbibliography[title={منابع و مراجع}]
\end{latin}

\begin{remark}
    برای به‌روزرسانی کتابنامه پس از افزودن استناد جدید،
    باید دستور \texttt{biber} را اجرا کنید:
    
    \begin{latin}
\begin{minted}{bash}
xelatex -shell-escape main.tex
biber main
xelatex -shell-escape main.tex
xelatex -shell-escape main.tex
\end{minted}
    \end{latin}
\end{remark}

%----------------------------------------
% Summary
%----------------------------------------
\chapter{خلاصه نتایج}

\section{وضعیت سیستم کتابنامه}

\begin{table}[h]
\centering
\caption{وضعیت قابلیت‌های کتابنامه}
\label{tab:bibliography-status}
\begin{tabular}{|c|C{6cm}|c|}
\hline
\textbf{قابلیت} & \textbf{توضیحات} & \textbf{وضعیت} \\
\hline
\lr{\textbackslash cite} & استناد ساده به منابع & \textcolor{matingreen}{✅} \\
\lr{\textbackslash parencite} & استناد پرانتزی & \textcolor{matingreen}{✅} \\
\lr{\textbackslash textcite} & استناد متنی & \textcolor{matingreen}{✅} \\
استناد چندگانه & چند منبع در یک استناد & \textcolor{matingreen}{✅} \\
\lr{\textbackslash printbibliography} & چاپ کتابنامه & \textcolor{matingreen}{✅} \\
شیوه عددی & \lr{numeric} با ترتیب استناد & \textcolor{matingreen}{✅} \\
عنوان فارسی & «منابع و مراجع» & \textcolor{matingreen}{✅} \\
فایل \lr{.bib} & مدیریت منابع خارجی & \textcolor{matingreen}{✅} \\
\hline
\end{tabular}
\end{table}

\section{نتیجه}

\textbf{سیستم کتابنامه \lr{MatinBook} نسخه ۱ با موفقیت آزمایش شد.}
تمامی دستورات استناد (\lr{cite}، \lr{parencite}، \lr{textcite})
و نمایش کتابنامه با \lr{biblatex} و \lr{biber}
به درستی کار می‌کنند. 🎉

\end{document}